%%%%%%%%%%%%%%%%%%%%%%%%%%%%%%%%%%%%%%%%%%%%%%%%%%%%%%%%%%%%%%%
% DBDBD 2026 - Official Template for Abstract Submission
% Filled in for the federated-fog scheduling paper.
% One page maximum INCLUDING references.
%%%%%%%%%%%%%%%%%%%%%%%%%%%%%%%%%%%%%%%%%%%%%%%%%%%%%%%%%%%%%%%
\documentclass[11pt]{article}

\usepackage{amssymb,amsmath,amsfonts,latexsym,graphicx}
\usepackage[left=2.5cm,right=2.5cm,top=1.5cm,bottom=2.5cm,footnotesep=0.5cm]{geometry}

\setlength{\oddsidemargin}{0.5cm}
\setlength{\textwidth}{15.5cm}
\setlength{\topmargin}{-1.5cm}
\setlength{\textheight}{22cm}

\pagestyle{empty}

\begin{document}

\begin{center}
    \Large{\textbf{Data-Aware Scheduling of Approximate DAG Workflows in Federated Fog Systems}}

    \vspace{0.5cm}

    \large{Gayani Gupta$\,^1$, Amisha Gupta$\,^2$, Mohsen Amini Salehi$\,^1$,
           Antonio Fern\'andez Anta$\,^3$, Sanjukta Bhowmick$\,^1$, and Jose Aguilar$\,^4$}

    \vspace{0.5cm}

    \normalsize

    $^1$ Dept.\ of Computer Science and Engineering, University of North Texas, USA\\
    \texttt{\{gayani.gupta, mohsen.aminisalehi, sanjukta.bhowmick\}@unt.edu}

    $^2$ University of California, Berkeley, USA

    $^3$ IMDEA Networks Institute, Madrid, Spain\\
    \texttt{antonio.fernandez@imdea.org}

    $^4$ Universidad de Los Andes, M\'erida, Venezuela\\
    \texttt{aguilar@ula.ve}

    \vspace{0.6cm}

\end{center}

% Text of the abstract (one page maximum, including references)
Data-intensive workflows such as disaster-response pipelines exchange large
intermediate results between dependent tasks, so where a task runs determines
whether its input arrives in time. We formulate this as a mixed-integer
program that jointly decides task placement across federated fog nodes and
the degree of approximate execution each task performs, subject to per-edge
data sizes, per-link bandwidths, node eligibility, and a workflow deadline.
We derive scheduling algorithms from the formulation: an exact MILP solver,
a randomized rounding of its linear relaxation, a data-aware greedy heuristic,
and a genetic algorithm. We evaluate them on synthetic DAGs of varying size,
deadline slack, federation size, link bandwidth, and task-failure rate. The
experiments show when each method earns its keep: under tight deadlines the
optimization-based methods lead, with rounding staying within a few percent of
the MILP, while the greedy method is near-optimal at the largest sizes at a
fraction of the solve time; a bandwidth-blind baseline quantifies what
data-aware placement is worth as links degrade. The work extends our earlier
conference study~\cite{gupta2025efficient} with the full formulation, its
linearization and proofs, and an analysis of scalability limits and task
failures.

\bibliographystyle{plain}
\bibliography{cas-refs}

\end{document}
